% Abhakseite „Das kann ich“ – Zone nur im Gesamt (\mitzone)
%% TODO Ich-kann-Sätze: Platzhalter wie in den Aufgabendateien
\begin{abhakseite}
\ifdefined\mitzone
\abhakgruppe{Kennst du schon}
\abhak{1}{Anteile, Prozentsätze und absolute Häufigkeiten ineinander umrechnen}
\abhak{2}{Zweistufige Zufallsexperimente und Pfadregeln (Baumdiagramm als Schwesterdarstellung, „Anteil vom Anteil“ als Produkt)}
\abhak{3}{Gegenereignis und Komplement (der Rest zu eins bzw. zur Gesamtzahl)}
\abhak{4}{Tabellen lesen und ausfüllen (Zeilen- und Spaltensummen)}
\abhak{5}{Lineare Gleichungen und Ungleichungen mit einem Parameter lösen (negative Werte ausschließen)}
\abhak{6}{Zweistufige Zufallsexperimente und Pfadregeln (Baumdiagramm als Schwesterdarstellung, „Anteil vom Anteil“ als Produkt) – Fehler finden}
\abhak{7}{Zweistufige Zufallsexperimente und Pfadregeln (Baumdiagramm als Schwesterdarstellung, „Anteil vom Anteil“ als Produkt) – selbst rechnen}
\fi
\abhakgruppe{Einheit 1 · Einheit 1}
\abhak{8}{Weder–noch, entweder–oder, oder?}
\abhak{9}{Die Tafel füllen}
\abhak{10}{Die Tafel füllen – Fehler finden, Begründen}
\abhakgruppe{Einheit 2 · Einheit 2}
\abhak{11}{Aus der Tafel rechnen}
\abhak{12}{Aus der Tafel rechnen – Fehler finden, Begründen}
\end{abhakseite}
